\documentclass[10pt,a4paper]{amsart}
\usepackage[margin=19mm]{geometry}
\usepackage{amsmath,amssymb,mathtools}
\usepackage[scr=rsfs,cal=euler]{mathalfa}
\usepackage{xfrac}
\usepackage[unicode,hidelinks,bookmarks=false]{hyperref}
\newcommand{\ie}{i.e.\ }
\newcommand{\cf}{cf.\ }
\newcommand{\resp}{respectively\ }
\newcommand{\Subsection}{\subsection}
\providecommand{\nobreakdash}{\nobreak\hbox{-}}
\setlength{\emergencystretch}{2em}
\multlinegap0pt
\usepackage{amssymb}
\usepackage{mathtools}
\usepackage{tikz}
\usepackage{float}
\usepackage{enumerate}
\usepackage[shortlabels]{enumitem}
\usepackage{csquotes}
\newtheorem{lemma}{Lemma}
\title{Eco-evolutionary dynamics of trait-structured, sexually reproducing interacting populations with time-dependent birth and mortality rates}
\author{Manh Hong Duong, Fabian Spill, Blaine van Rensburg}
\date{Source: arXiv:2501.07379v3 (CC BY 4.0)}
\begin{document}
\maketitle

\noindent\emph{Source and licence.} Eco-evolutionary dynamics of trait-structured, sexually reproducing interacting populations with time-dependent birth and mortality rates. Version arXiv:2501.07379v3 (CC BY 4.0), distributed by arXiv under the Creative Commons Attribution 4.0 International licence, http://creativecommons.org/licenses/by/4.0/, as recorded in the arXiv licence metadata for this exact version and shown on the abstract page https://arxiv.org/abs/2501.07379v3. Full licence notice: \texttt{licenses/arxiv-2501.07379-CC-BY-4.0.txt}.\par\smallskip

\noindent\textit{Editorial scope.} Counted unit: the article's lemma lma:KPPBound (source0.tex lines 842--856). The article's complete original proof of it is delivered with it (source0.tex lines 858--895, one \texttt{proof} environment of 38 lines). No mathematics is written, repaired or re-derived: the statement and the proof are the article's own lines, in the article's own notation, and no retained line is substituted or re-flowed.  No further source-local context is needed: every cross-reference of every retained line resolves inside the two delivered spans.  Nothing was omitted from any retained span, so the package carries no omission record.  No retained line contains a citation, so the wrapper carries no bibliography and no bibliography entry.  No retained line contains a figure environment, an image-inclusion command or drawing code, so no image is delivered and the package declares no asset dependency.  Editorial formatting only, in addition: an AMS \texttt{amsart} preamble that restates the article's own declarations and macros used by the retained lines, namely the environments \texttt{lemma} on separate counters, together with the standard packages those lines need.  The retained environments print in this document as Lemma 1 in order of appearance, whereas the article prints the same items with the numbering of its own full document.  The complete native source of the article as distributed in the arXiv e-print for this exact version is bundled unchanged as \texttt{sources/171548/source0.tex}.
\begin{lemma}\label{lma:KPPBound}

Suppose that $y_\varepsilon(t)$ solves 
\begin{equation}\label{eqn:GenKPP}
    \begin{cases}
\varepsilon^2\frac{dy_\varepsilon}{dt}=\alpha_\varepsilon(t)(H_\varepsilon(t)-y_\varepsilon)y_\varepsilon,\\
y_\varepsilon(0)=H_\varepsilon(0)+O(\varepsilon),
\end{cases}
\end{equation}
where there exit constants $\alpha_L,\alpha_U,H_L,H_U$ such that $0<\alpha_L<\alpha_\varepsilon(t)<\alpha_U$, $0<H_L<H_\varepsilon(t)<H_U$, $H_\varepsilon'(t)<C_H$, and $\alpha'_\varepsilon(t)<C_\alpha$. Then 
\begin{itemize}
    \item[(a)] $|y_\varepsilon(t)-H_\varepsilon(t)|\leq{}C\varepsilon$, where $C$ is a constant depending on $C_H$, $H_L$, $H_U$ and initial condition.
    \item[(b)] $\dot{y}_\varepsilon(t)\leq{}C_1+\varepsilon^{-1}C_2\exp^{-\frac{H_L}{2\varepsilon^2}t}.$
\end{itemize}
\end{lemma}

\begin{proof}[Proof of Lemma \ref{lma:KPPBound}]
Letting $J_\varepsilon(t):=\frac{1}{y_\varepsilon(t)}$ we find that 
\[J_\varepsilon(t)=J_\varepsilon(0)e^{-\frac{\int_{0}^{t}\alpha_\varepsilon(s)H_\varepsilon(s)ds}{\varepsilon^2}}+\varepsilon^{-2}\int_{0}^{t}\alpha_{\varepsilon}(s)e^{-\frac{1}{\varepsilon^2}\int_{s}^{t}H_\varepsilon(\tau)\alpha_\varepsilon(\tau)d\tau}ds.\]
We focus for now on the latter term. Let $\beta\in(1,2)$, and suppose $t<\varepsilon^\beta$. Then for $0<s<t$ we have
\[|H_\varepsilon(s)-H_\varepsilon(t)|\leq{}C_H{}\varepsilon^\beta,\]
 due to the $\varepsilon$ independent bounds on $H_\varepsilon(t)$. We obtain, using also the uniform lower bounds on $\alpha_\varepsilon(t)$ and $H_\varepsilon(t)$,
\[J_\varepsilon(t)\leq{}J_\varepsilon
(0)e^{-\frac{H_\varepsilon(t)-C_H\varepsilon^\beta}{\varepsilon^2}\int_{0}^{t}\alpha_\varepsilon(\tau)d\tau}+\varepsilon^{-2}\int_{0}^{t}\alpha_\varepsilon(s)e^{-\frac{H_\varepsilon(t)-C_H\varepsilon^\beta}{\varepsilon^2}\int_{s}^{t}\alpha_\varepsilon(\tau)d\tau}ds.\]
A change of variables allows us to simplify the RHS to give 
\[J_\varepsilon(t)\leq{}J_\varepsilon
(0)e^{-\frac{H_\varepsilon(t)-C_H\varepsilon^\beta}{\varepsilon^2}\int_{0}^{t}\alpha_\varepsilon(\tau)d\tau}+\frac{1}{H_\varepsilon(t)-C_H\varepsilon^\beta}\left(1-e^{-\frac{H_\varepsilon(t)-C_H\varepsilon^\beta}{\varepsilon^2}\int_{0}^{t}\alpha_\varepsilon(\tau)d\tau}\right).\]
A similar set of computations provides a lower bound
\[J_\varepsilon(t)\geq{}J_\varepsilon
(0)e^{-\frac{H_\varepsilon(t)+C_H\varepsilon^\beta}{\varepsilon^2}\int_{0}^{t}\alpha_\varepsilon(\tau)d\tau}+\frac{1}{H_\varepsilon(t)+C_H\varepsilon^\beta}\left(1-e^{-\frac{H_\varepsilon(t)+C_H\varepsilon^\beta}{\varepsilon^2}\int_{0}^{t}\alpha_\varepsilon(\tau)d\tau}\right).\]
From this we obtain, for $t\in(0,\varepsilon^\beta)$, \[J_\varepsilon\in\left(\min\left\{J_\varepsilon(0),\frac{1}{H_\varepsilon(t)+C_H\varepsilon^\beta}\right\},\max\left\{J_\varepsilon(0),\frac{1}{H_\varepsilon(t)-C_H\varepsilon^\beta}\right\}\right).\]
We now use the fact that $y_\varepsilon(0)=H_\varepsilon(0)+O(\varepsilon)=H_\varepsilon(t)+O(\varepsilon)$, to deduce that $y_\varepsilon(t)=H_\varepsilon(t)+O(\varepsilon)$ for $t\in[0,\varepsilon^\beta].$


If $t>\varepsilon^\beta$ instead, we find that the first term is exponentially small, and can estimate the latter similarly to above. Overall this shows $y_\varepsilon(t)=H_\varepsilon(t)+O(\varepsilon)$. This completes part $(a)$. 

We now show part $(b)$. First, we differentiate \eqref{eqn:GenKPP} and set $\nu_\varepsilon=\dot{y}_\varepsilon$.  We obtain 
\begin{align*}
\varepsilon^{2}\dot{\nu}_\varepsilon&=(H_\varepsilon(t)-y_\varepsilon)\nu_\varepsilon-(H_\varepsilon'(t)-\nu_\varepsilon)y_\varepsilon\\
&=-H_\varepsilon'(t)y_\varepsilon-(2y_\varepsilon-H_\varepsilon(t))\nu_\varepsilon.
\end{align*}
Using part (a), we find that $2y_\varepsilon-H_\varepsilon=H_\varepsilon(t)+O(\varepsilon)$ and so
\[\varepsilon^{2}\dot{v}_\varepsilon\leq{}2C_HH_U-\frac{H_L}{2}\nu_\varepsilon,\]
when $\nu_\varepsilon>0$ and
\[\varepsilon^{2}\dot{v}_\varepsilon\geq{}-2C_HH_U-\frac{H_L}{2}\nu_\varepsilon,\]
when $\nu_\varepsilon<0$. We therefore have, for $\overline{\nu}_\varepsilon=|\nu_\varepsilon|$,
\[\varepsilon^2\overline{\nu}_\varepsilon\leq{}2C_HH_U-\frac{H_L}{2}\overline{\nu}_\varepsilon,\]
from which we derive
\[
|\nu_\varepsilon|\leq{}2C_H\rho_U\left(1-e^{-\frac{H_L}{2\varepsilon^2}t}\right)+|\nu_\varepsilon(0)|e^{-\frac{H_L}{2\varepsilon^2}t}.
\]
Altogether, this gives $|\nu_\varepsilon|\leq{}C_1+C_2\varepsilon^{-1}e^{-\frac{H_L}{2\varepsilon^2}t}$, proving part (b). 

\end{proof}

\end{document}
